\section{Hyperbolicity and two cone slices}
\label{sec:hyperbolic}

We first prove that the polynomial in \eqref{eq:polynomial} is
hyperbolic.  We then identify the two slices of its closed cone that
will constrain any representing pencil.

For $u=(u_0,u')$ and $v=(v_0,v')$ in $\R^4$, the block formula
\eqref{eq:phi} gives
\begin{equation}
 u^\top\Phi_y(vv^\top)u
 = (v_0u'-u_0v')^\top Q(y)(v_0u'-u_0v')
 = b(v_0u'-u_0v',y).
 \label{eq:rank-one}
\end{equation}
Indeed, the three terms on expanding the middle expression are
$v_0^2u'^\top Q(y)u'$, $-2u_0v_0u'^\top Q(y)v'$, and
$u_0^2v'^\top Q(y)v'$, exactly the terms from the block formula.
Lemma~\ref{lem:form} therefore implies that
$\Phi_y(vv^\top)\succeq0$.  Every positive semidefinite real symmetric
matrix is a nonnegative linear combination of such rank-one matrices.
Thus $\Phi_y$ is a positive linear map: it preserves positive
semidefiniteness, and hence the order $\preceq$.
The same formula also gives
\begin{equation}
 \Phi_{s y}=s^2\Phi_y\qquad(s\in\R),
 \label{eq:phi-scaling}
\end{equation}
so in particular $\Phi_0=0$ and $\Phi_{-y}=\Phi_y$.
This rank-one verification is the real-symmetric positive-map
correspondence for biquadratic forms~\cite[Section~3]{ha-2013}
applied to the explicit map~\eqref{eq:phi}.

\begin{proposition}\label{prop:hyperbolic}
The polynomial $p$ in \eqref{eq:polynomial} is homogeneous of
degree $20$ on $\Sym^4\times\Sym^4\times\R^3$, a real vector
space of dimension $23$.  It satisfies $p(e)=1$ and is hyperbolic
with respect to $e=(I_4,I_4,0)$.
\end{proposition}

\begin{proof}
Each entry of $(\det X)Z-\Phi_y(\operatorname{adj}X)$ is homogeneous
of degree $5$: the adjugate has degree $3$, and $\Phi_y$ is quadratic
in $y$.  Taking its $4\times4$ determinant gives a homogeneous
polynomial of degree $20$.  Since $\Phi_0=0$, evaluation at $e$
gives $p(e)=1$, so this polynomial is nonzero.  The dimension is
$10+10+3=23$.

Extend $\Phi_y$ complex linearly to complex symmetric matrices.
For invertible $X$, real or complex, linearity and
$\operatorname{adj}X=(\det X)X^{-1}$ give
\begin{equation}
 p(X,Z,y)=(\det X)^4\det\bigl(Z-\Phi_y(X^{-1})\bigr).
 \label{eq:inverse-formula}
\end{equation}
The adjugate formula remains the definition at singular $X$.

Fix real symmetric $X,Z$ and real $y$, and let $t=a+i\eta$ with
$\eta>0$.  Orthogonal diagonalization of $X$ gives
\[
 \operatorname{Im}(tI_4-X)^{-1}
 =-\eta\bigl((aI_4-X)^2+\eta^2I_4\bigr)^{-1}\prec0.
\]
For a complex symmetric matrix, the entrywise imaginary part equals
the Hermitian imaginary part $(A-A^*)/(2i)$ and is real symmetric.
Complex linearity and positivity of $\Phi_{-y}$ consequently show
that
\[
 H=tI_4-Z-\Phi_{-y}\bigl((tI_4-X)^{-1}\bigr)
 \quad\text{satisfies}\quad
 \operatorname{Im}H\succeq\eta I_4\succ0.
\]
Such an $H$ is invertible: if $Hw=0$ for a nonzero complex vector
$w$, then $0=\operatorname{Im}(w^*Hw)=w^*(\operatorname{Im}H)w>0$.
Also $tI_4-X$ is invertible.  Equation~\eqref{eq:inverse-formula}
therefore implies $p(te-(X,Z,y))\ne0$.
The coefficients of this polynomial in $t$ are real, so conjugation
also excludes roots in the lower half-plane.  Finally, its leading
coefficient is $p(e)=1$ by homogeneity; it has degree $20$ for every
real input.  All its roots are therefore real, as required.
\end{proof}

\begin{proposition}\label{prop:slices}
Let $K=\Lambda_+(p,e)$.  For real symmetric $X,Z$ and $y\in\R^3$,
\begin{align}
 X\succ0:\qquad
 (X,Z,y)\in K
 &\ \Longleftrightarrow\ 
 Z\succeq\Phi_y(X^{-1}),
 \label{eq:slice-positive}\\
 (X,Z,0)\in K
 &\ \Longleftrightarrow\ 
 X\succeq0\ \text{and}\ Z\succeq0.
 \label{eq:slice-zero}
\end{align}
Moreover, $e$ lies in the interior of $K$ in the full ambient space.
\end{proposition}

\begin{proof}
If $\lambda_1(x),\ldots,\lambda_{20}(x)$ are the roots of
$t\mapsto p(te-x)$, homogeneity gives
\[
 p(x+te)=\prod_{j=1}^{20}(t+\lambda_j(x)).
\]
Thus $x\in K$ if and only if $p(x+te)\ne0$ for every real $t>0$.
The strict inequality on $t$ allows zero roots and hence includes
the boundary of $K$.

Suppose $X\succ0$ and set, for $t\ge0$,
\[
 F(t)=Z+tI_4-\Phi_y\bigl((X+tI_4)^{-1}\bigr).
\]
If $F(0)\succeq0$, order preservation yields
\[
 F(t)=F(0)+tI_4+
 \Phi_y\bigl(X^{-1}-(X+tI_4)^{-1}\bigr)
 \succeq tI_4\succ0\qquad(t>0).
\]
Equation~\eqref{eq:inverse-formula} then shows that there is no
positive root, so $(X,Z,y)\in K$.

Conversely, if $F(0)\not\succeq0$, its least eigenvalue is negative.
For $t>0$, positivity gives
\[
 0\preceq\Phi_y\bigl((X+tI_4)^{-1}\bigr)
 \preceq t^{-1}\Phi_y(I_4).
\]
Consequently $F(t)\succ0$ for all sufficiently large $t$.
Continuity of its least eigenvalue gives a $t>0$ at which $F(t)$
is singular.  Since $X+tI_4\succ0$,
Equation~\eqref{eq:inverse-formula} now gives a positive root of
$p((X,Z,y)+te)$.  This proves \eqref{eq:slice-positive} in both
directions, including equality in the matrix inequality.

When $y=0$, the polynomial definition directly gives
\[
 p(te-(X,Z,0))=\det(tI_4-X)^4\det(tI_4-Z),
\]
also for singular $X$ and $Z$.  Its roots are the eigenvalues of $X$,
each repeated four times, and the eigenvalues of $Z$.
They are all nonnegative exactly when both matrices are positive
semidefinite.  This proves \eqref{eq:slice-zero}.

Finally, $X\succ0$ and $Z-\Phi_y(X^{-1})\succ0$ define an open
subset of the full ambient space.  It contains $e$ and, by
\eqref{eq:slice-positive}, is contained in $K$.  Thus $e$ is an
interior point.
\end{proof}
